\documentclass[../main.tex]{subfiles}

\begin{document}

\ifSubfilesClassLoaded{

    %\tableofcontents
    
    \setcounter{chapter}{6}
}{}

\chapter{ HW7 }
 代靖涵 25120222201319

\begin{exercise}{}{}
\begin{enumerate}
    \item[(1)] 求证对任意 $X, Y \in \Gamma(TM)$, 有
    \[
        [X,Y] := XY - YX \in \Gamma(TM).
    \]
    \item[(2)] 证明 Jacobi 恒等式
    \[
        [[X,Y],Z] + [[Z,X],Y] + [[Y,Z],X] = 0, \quad \forall X,Y,Z \in \Gamma(TM).
    \]
\end{enumerate}
\end{exercise}
\begin{proof}
    \begin{enumerate}
        \item 任取 \(  f,g \in C^{\infty}\left( M \right)   \),  \[
    \begin{aligned}
    [X,Y]\left( fg \right)&= X\left( Y\left( fg \right)  \right)-Y\left( X\left( fg \right)  \right)\\ 
     &=  X\left( f Yg+ gYf \right)- Y\left( fXg+ gXf \right)\\ 
      &= f \left( XYg \right) + \left( Xf \right)\left( Yg \right)+  \left( Xg \right)\left( Yf \right)+ g\left( XYf \right)\\ 
       &- \left( Yf \right)\left( Xg \right)- f\left( YXg \right)- \left( Yg \right)\left( Xf \right)- g\left( YXf \right)\\ 
        &= f\left( XYg \right)+ g\left( XYf \right)- f\left( YXg \right)-            g\left( YXf \right)\\ 
         &= f\left(  \left[ X,Y \right]g   \right)+ g\left( \left[ X,Y \right]f  \right)  
    \end{aligned}
    \]这表明 \(  \left[ X,Y \right]   \)满足Lebniz率. 此外对于任意的 \(  k \in \mathbb{R}   \),   \[
   \begin{aligned}
    [X,Y]\left( kf + g\right)&= XY\left( kf + g\right)-YX\left( kf + g\right)\\ 
     &= k\left( XY-YX \right)f +\left( XY-YX \right)g  = k\left[ X,Y \right]\left( f \right)   + \left[ X,Y \right]g  
   \end{aligned}  
    \]故 \(  \left[ X,Y \right]   \)是\(  C^{\infty}\left( M \right)   \)上的一个导子,  \(  \left[ X,Y \right] \in  \Gamma \left( TM \right)    \).   
    \item  \[
    \begin{aligned}
    \left[ \left[ X,Y \right],Z  \right]&= \left[ XY-YX,Z \right]= \left( XY-YX \right)Z-Z\left( XY-YX \right)\\ 
     &=  XYZ-YXZ-ZXY+ ZYX     
    \end{aligned}
    \]由对称性, 我们得到 \[
    \left[ \left[ Z,X \right],Y  \right]= ZXY- XZY- YZX+ YXZ
    \] \[
    \left[ \left[ Y,Z \right],X  \right]=YZX- ZYX-XYZ+ XZY
    \]将以上三式相加, 发现右侧12项两两抵消, 因此 \[
    \left[ \left[ X,Y \right],Z  \right]+ \left[ \left[ Z,X \right],Y  \right]+ \left[ \left[ Y,Z \right],X  \right]   = 0
    \]
    \end{enumerate}
    
\end{proof}
\begin{exercise}{}{}
定义沿着闭集上的光滑向量场: 对任意闭子集 $A \subset M$, 称 $Y$ 是沿着 $A$ 定义的一个光滑向量场, 若 $\forall p \in A, Y_p \in T_p M$, 且存在 $p$ 的某个领域 $U$ 和 $U$ 上定义的光滑向量场 $\tilde{Y}$, 使得 $\tilde{Y}|_{A \cap U} = Y$. 一个向量场的支集定义为 $\operatorname{supp}(X) = \overline{\{x \in M | X(x) \neq 0\}}$. 求证给定任意闭集 $A$, 存在包含 $A$ 的开集 $A \subset U$, 和向量场 $\tilde{Y} \in \Gamma(TM)$ 使得 $\tilde{Y}|_A = Y$ 并且 $\operatorname{supp}(\tilde{Y}) \subset U$. 特别地, 若 $A = \{p\}$, 我们得到单个向量的延拓定理.
\end{exercise}

\begin{proof}
    设 \(  Y  \)是沿着 \(  A  \)的光滑向量场. 对于任意的 \(  p \in A \), 设 \(  U_{p}  \)是\(  p  \)的邻域,  \(  \tilde{Y}^{\left( p \right) }  \)是 \(  U _{p} \)上的光滑向量场,        使得\(  \left. \tilde{Y}^{\left( p \right) } \right|_{A\cap U_{p}}= Y  \) . \(  \left\{ U_{p} \right\}_{p \in A}\cup \left\{ M\setminus A \right\}  \)构成 \( M  \)的一个开覆盖. 取从属于 \(  \left\{ U_{p} \right\} \cup \left\{ M\setminus A \right\} \)的光滑单位分解 \(  \left\{ \psi _{p} \right\} \cup \left\{  \varphi  \right\} \), 使得 \(  \operatorname{supp}\psi _{p}\subseteq U_{p}  ,\forall p \in A\), \(  \operatorname{supp} \varphi \subseteq M\setminus A  \). 令 \(  U= \bigcup _{p \in A}U_{p}  \) , 对于每个 \(  p  \), 定义向量场 \[
    Z_{p}= \begin{cases} \psi _{p}\left( q \right) \tilde{Y}^{\left( p \right) }_{q},& q \in U_{p}\\ 
     0,& q \not \in  \operatorname{supp}\psi _{p}\end{cases} 
    \]   再定义 \[
    \tilde{Y}= \sum _{p \in A} Z_{p}
    \] 由于单位分解是局部有限的, 因此 上述表达式在 \(  U  \)上任一点都是一个有限和,  从而 \(  \tilde{Y}  \)是良定义的.  为了说明 \(  \tilde{Y}|_{A}= Y  \), 任取  \(  q \in A  \), 设 \(  V  \)是 \(  q  \)的一个邻域, 使得 \(  \left\{ \psi _{p} \right\}  \)中支撑集与 \(  V  \)相交的 元素的只有   \(  \psi_{p_1},\cdots ,\psi _{p_{k}}  \)    和 \(  \varphi   \) 从而\[
     1= \sum _{p \in A}\psi _{p}\left( q \right)+  \varphi \left( q \right)  = \sum _{i= 1}^{k}\psi _{p_{i}}\left( q \right) +  \varphi \left( q \right) 
    \] 但是 \[
    q \in A,\quad  \operatorname{supp} \varphi \subseteq M\setminus A\implies q\not \in \operatorname{supp} \varphi \implies  \varphi \left( q \right)= 0 
    \]因此 \[
    \sum _{i= 1}^{k}\psi _{p_{i}}\left( q \right)= 1 
    \]\[
    \tilde{Y}_{q}= \sum _{i= 1}^{k}\psi _{p_{i}}\left( q \right) \tilde{Y}_{q}^{\left( p_{i} \right) } = \sum _{i= 1}^{k}\psi _{p_{i}}\left( q \right) Y_{q}= Y_{q}
    \]即 \[
    \tilde{Y}|_{A}= Y
    \]接下来, 为了说明 \(  \tilde{Y}  \)是 \(  M  \)上的光滑向量场, 我们需要说明它是一个 \(  C^{\infty}\left( M \right)   \)-导子, 为此, 任取 \(  f,g \in C^{\infty}\left( M \right)   \)     \[
   \begin{aligned}
    \tilde{Y}_{q}\left( fg \right)&= \left( \sum _{i= 1}^{k}\psi _{p_{i}}\left( q \right)  \tilde{Y}_{q}^{\left( p_{i} \right) }\right)\left( fg \right) \\ 
     &= \sum _{i= 1}^{k}\psi _{p_{i}}\left( q \right)  \tilde{Y}_{q}^{\left( p _{i}\right) }\left( fg \right)\\ 
      &= \sum _{i= 1}^{k}\psi _{p_{i}}\left( q \right)  \left( f \tilde{Y}_{q}^{\left( p_{i} \right) }g+  g \tilde{Y}_{q}^{\left( p_{i} \right) }f \right)= f \left( \sum _{i= 1}^{k}\psi _{p_{i}}\left( q \right)  Y_{q}^{\left( p_{i} \right) } \right) g+  g \left( \sum _{i= 1}^{k}\psi_{p_{i}}\left( q \right) Y_{q}^{\left( p_{i} \right) } \right)f \\ 
       &= f \tilde{Y}_{q}g+  g \tilde{Y}_{q}f
   \end{aligned}   
    \]此外, 由于对于任意的 \(  q \not \in U  \),  \(  \tilde{Y}_{q}= 0  \), 显然也满足上述关系.  因此 \[
    \tilde{Y} \left( fg \right)=  f \tilde{Y}g+ g \tilde{Y} f 
    \]并且对于任意的 \(  l\in  \mathbb{R}   \), 可以类似地说明 \[
    \tilde{Y} \left( lf+ g \right)= l \tilde{Y}f+ \tilde{Y}g 
    \]   从而 \(  \tilde{Y}  \)是 \(  M  \)上的光滑向量场.  根据\(  \tilde{Y}  \)的定义,  \[
    \operatorname{supp}\left( \tilde{Y} \right)\subseteq   \overline{\bigcup _{p \in A} \operatorname{supp}\psi _{p}}= \bigcup _{p \in A} \operatorname{supp}\psi _{p}\subseteq \bigcup _{p \in A}U_{p}= U 
    \]   

    上述构造的集合 \(  U  \)和向量场 \(  \tilde{Y}  \)即为所需. 特别地, 若 \(  A= \left\{ p \right\}  \), 我们得到单个向量场的延拓定理.   
\end{proof}
\begin{exercise}{}{}
\begin{enumerate}
    \item[a).] 显式构造 $S^{2n+1}$ 上处处非零的切向量场, 并说明 $S^{2n}$ 上能否构造;
    \item[b).] 构造 $\mathbb{T}^n$ 上处处非零的切向量场;
    \item[c).] 构造 $S^2$ 上仅有两个零点的向量场;
    \item[d).] 构造 $S^2$ 上仅有一个零点的向量场.
\end{enumerate}
\end{exercise}
\begin{solution}
    \begin{enumerate}
        \item[a).] 考虑标准的嵌入 \(  S^{2n+ 1}\hookrightarrow \mathbb{R} ^{2n+ 2}  \), 并考虑它在欧式空间上的标准坐标表示. 定义 \(  S^{2n+ 1}  \)上的一个向量场, \[
        X= \sum _{i= 1}^{n+ 1}\left(  x^{2i-1}  \frac{\partial }{\partial x^{2i}}-x^{2i} \frac{\partial }{\partial x^{2i-1}} \right) 
        \]  为了说明 \(  X \in  \Gamma \left(T S^{2n+ 1} \right)   \) , 只需要注意到在 \(  S^{2n+ 1}  \)任一点处, \(  X  \)在该点的切向量等同于 \(  \mathbb{R} ^{2n+ 2}  \)上的向量 \(  \left( -x^{2},x^{1},\cdots ,-x^{2n},x^{2n+ 1} \right)   \), 并且它与位置向量的内积     \[
        \left( x^{1},\cdots ,x^{2n+ 2} \right)\cdot  \left(- x^{2},x^{1},\cdots ,- x^{2n+ 2}, x^{2n+ 1} \right)  = 0
        \]这表明 \(  X  \)在任一点处的取值确实位于 \(  TS^{2n+ 1}  \)  ,  并且由于 \(  X  \)的分量函数是光滑的, 且处处不全为零(否则位置向量为零, 不属于球面). 因此 \(  X  \)是 \(  S^{2n+ 1}  \)上的一个处处非零的光滑向量场.   



        不存在 \(  S^{2n}  \)上非退化的光滑切向量场. 事实上, 若存在这样的向量场 \(  X  \in  \Gamma \left( TS^{2n} \right) \). 考虑 \(  S^{2n}  \)到 \(  \mathbb{R} ^{2n+ 1}  \)的典范嵌入. 此时,对于任意一点 \(  p\in  S^{2n} \), 我们有切空间的典范嵌入和典范同构 \(  T_{p}S^{2n}\hookrightarrow T_{p}\mathbb{R} ^{2n+ 1}\simeq \mathbb{R} ^{2n+ 1}  \) ,  此时 \(  T_{p}S^{2n}\subseteq \mathbb{R} ^{2n+ 1}  \)被赋予了标准的坐标和内积. 由于 \(  X  \)是非退化的,  \(  \left<X_{p},X_{p} \right>> 0  \), 从而可以对 \(  X  \)做单位化, 得到向量场 \(  Y  \) :     \[
        Y_{p}= \frac{X_{p} }{\sqrt{\left<X_{p}, X_{p} \right>} } 
        \]由于内积是光滑的映射, \(  Y  \)是光滑的向量场 . 此时, 将 \(  Y_{p}  \)用 \(  \mathbb{R} ^{2n+ 1}  \)上的标准坐标表示, 它   是 \(  S^{2n}  \)上的向量. 由于 \(  \left\| Y_{p} \right\|= 1, \forall p \in S^{2n}  \), 我们给出 一个光滑映射\(  S^{2n}\to S^{2n}  \) \[
        p\mapsto  Y_{p}
        \]  进而, 我们可以构造同伦映射 \(  H: S^{2n}\times \mathbb{I}\to S^{2n}  \) \[
        H\left( p,t \right)=  \cos  \left(  \pi t \right) p   + \sin \left(  \pi t \right)  Y_{p} 
        \] 由于 \(  Y_{p}\in T_{p}S^{2n}  \), \(  Y_{p}  \)与 \(  p  \)正交, 由 \(  Y_{p},p  \)均为单位向量  因此    \[
        \left<H\left( p,t \right),H\left( p,t \right)   \right>= 1
        \]从而 \(  H\left( p,t \right) \in S^{2n}, \forall \left( p,t \right)\in S^{2n}\times \mathbb{I}    \) . 因此 \(  H  \)是良定义的. 且满足 \[
        H\left( p,0 \right)= p,\quad H\left( p,1 \right)= -p  
        \] 从而 \(  H  \)给出了单位映射\(  \operatorname{Id}_{S^{2n}}  \) 和对径映射 \(  A  \) 的同伦 \(  H: \operatorname{Id}_{S^{2n}}\simeq  A \)  . 但是对于 \(  A: S^{2n}\to S^{2n}  \),   \(  \operatorname{deg}A  = \left( -1 \right)^{2n+ 1}= -1 \), 而 \(  \operatorname{deg} \operatorname{Id}_{S^{2n}}= 1  \), 又映射的 \(  \operatorname{deg}  \)是同伦不变的, 矛盾. 因此不存在 \(  S^{2n}  \)上无处退化的光滑向量场.  

 
        
        \item[b).] \(  \mathbb{T}^{n}  \)微分同胚于 \(  n  \)个圆的乘积流形 \[
        S^{1}\times S^{1}\times \cdots \times S^{1}
        \]  我们知道 \(  S^{1}\subseteq \mathbb{R} ^{2}  \)上存在一个无处退化的光滑切向量 \[
        V_{\left( x,y \right) }= -y\frac{\partial }{\partial x}+ x\frac{\partial }{\partial y}
        \] 考虑切丛的同构 \[
        T \mathbb{T}^{n}\simeq TS^{1}\oplus TS^{1}\oplus \cdots \oplus TS^{1}
        \]对于任意的 \(  p= \left( p_1,p_2,\cdots ,p_{n} \right)\in S^{1}\times S^{1}\times \cdots S^{1}   \),  我们可以定义切向量场 \[
        X_{p}= \left( V_{p_1},0_{p_2},\cdots ,0_{p_{n}} \right) \in TS^{1}\oplus \cdots \oplus TS^{1}\simeq T \mathbb{T}^{n}
        \]则 \(  X  \)是\(  \mathbb{T}^{n}  \)上无处退化的光滑切向量场.  
       \item [c).]设 \(  \left( x_1,x_2,x_3 \right)   \)是\(  \mathbb{R} ^{3}  \)上的坐标, 将 \(  S^{2}  \)上点的位置和切向量均用标准欧式坐标表示,考虑 \(  S^{2}  \)上的向量场.    \[
        X_{\left( x_1,x_2,x_3 \right) }= -x_2\frac{\partial }{\partial x^{1}}+ x_1\frac{\partial }{\partial x^{2}}
        \]任意位置向量 \(  \left( x_1,x_2,x_3 \right)   \)与切向量 \(  X_{\left( x_1,x_2,x_3 \right) }= \left( -x_2,x_1,0 \right)   \)    的内积为 \[
        \left( x_1,x_2,x_3 \right)\cdot \left( -x_2,x_1,0 \right)= 0  
        \]因此 \(  X  \)是 \(  S^{2}  \)的切向量场.   \(  X= 0  \)当且仅当 \(  x_2= x_1= 0  \), 此时 \(  x_{3}= \pm  1 \), 因此\(  X  \)仅有两个零点\(  \left( 0,0,1 \right) , \left( 0,0,-1 \right)   \). 
        \item[d).]   沿用c).中的坐标表示, 考虑 \(  S^{2}  \)上的向量场.    \[
        X_{\left( x_1,x_2,x_3 \right) }= \left( x_1^{2}-x_3-1 \right)\frac{\partial }{\partial x_1}+ x_1x_2\frac{\partial }{\partial x_2}+ x_1\left( 1+ x_3 \right)\frac{\partial }{\partial x_3}  
        \]任意位置向量 \(  \left( x_1,x_2,x_3 \right)   \)与切向量 \(  X_{\left( x_1,x_2,x_3 \right) }= (x_1^2 - x_3 - 1, x_1x_2, x_1(1+x_3))  \)    的内积为 \[
       \begin{aligned}
        &\left( x_1,x_2,x_3 \right)\cdot (x_1^2 - x_3 - 1, x_1x_2, x_1(1+x_3)) \\ 
         &= x_1^{3}-x_1x_3-x_1+ x_1x_2^{2}+ x_1x_3+ x_1x_3^{2}\\ 
          &= x_1^{3}-x_1+  x_1\left( x_2^{2}+ x_3^{2} \right)\\ 
           &= x_1^{3}-x_1+ x_1\left( 1-x_1^{2} \right)= 0  
       \end{aligned} 
        \]因此 \(  X  \)是 \(  S^{2}  \)的切向量场.  若 \(  X= 0  \), 则 \(  x_1\left( 1+ x_3 \right)= 0   \), 从而 \(  x_1= 0  \)或 \(  x_3= -1  \), 但若 \(  x_1= 0  \), \(  x_1^{2}-x_3-1= 0  \)又会导致 \(  x_3= -1  \). 因此 使得 \(  X= 0  \)的点只有 \(  \left( 0,0,-1 \right)   \).        
    \end{enumerate}
    
\end{solution}

\begin{exercise}{}{}
设 $M$ 是欧氏平面 $\mathbb{R}^2$ 的开子流形, 坐标皆为正数. 定义映射
\[
F: M \to M, \quad (x, y) \mapsto (xy, \frac{y}{x}).
\]
求证
\begin{enumerate}
    \item $F$ 是一个微分同胚:
    \item 定义两个向量场,
    \[
    X = f(x)\frac{\partial}{\partial x} + g(y)\frac{\partial}{\partial y}, \quad Y = h(x, y)\frac{\partial}{\partial x},
    \]
    计算 $dF(X), dF(Y)$, 其中 $f, g, h$ 都是 $M$ 上的光滑函数:
    \item 计算 $[dF(X),dF(Y)]$ 并比较 $dF[X,Y]$. 这是一个普遍的规律吗, 给出你的理由.
\end{enumerate}
\end{exercise}

\begin{proof}
    \begin{enumerate}
        \item 由于坐标皆为正数, 故 \(  \left( x,y \right)\mapsto  xy   \)和 \(  \left( x,y \right)\mapsto \frac{y }{x }    \)均为 \(  M  \)上的光滑映射, 从而 \(  F  \)是光滑的. 此外,  \(  F  \)有逆映射  \(  F^{-1} :M\to M  \) ,\[
        F^{-1} \left( u,v \right)= \left( \sqrt{\frac{u }{v } }, \sqrt{uv} \right)  
        \]   也是一个光滑映射.  因此 \(  F  \)是一个微分同胚
        \item  在标准的欧式坐标下,  \(  \,\mathrm{d} F  \)的Jaocbi矩阵为 \[
        JF_{\left( x,y \right) }= \begin{pmatrix} 
            y& x\\ 
             -\frac{y }{x^{2} }& \frac{1 }{x }   
        \end{pmatrix} 
        \] \(  X _{\left( x,y \right) } \)的坐标为 \(  \left( f\left( x \right),g\left( y \right)   \right)^{\top}   \), 从而 \(  \,\mathrm{d} F\left( X \right)_{\left( x,y \right) }   \)的坐标为 \[
        JF_{\left( x,y \right) } \begin{pmatrix} 
            f\left( x \right)\\ 
             g\left( y \right)   
        \end{pmatrix}= \begin{pmatrix} 
              yf\left( x \right)+ xg\left( y \right)\\ 
               -\frac{yf\left( x \right)  }{x^{2} }+ \frac{g\left( y \right)  }{x }    
        \end{pmatrix}  
        \]   于是 \[
        \,\mathrm{d}_{\left( x,y \right) } F\left( X \right)= \left( yf\left( x \right)+ xg\left( y \right)   \right)\frac{\partial }{\partial u}+ \left( -\frac{yf\left( x \right)  }{x^{2} }+ \frac{g\left( y \right)  }{x }   \right)\frac{\partial }{\partial v}   
        \] 其中 \(  x= \sqrt{\frac{u }{v } }, y= \sqrt{uv}  \) 
        

         \[
         JF_{\left( x,y \right) }\begin{pmatrix} 
             h\left( x,y \right)\\ 
              0  
         \end{pmatrix}= \begin{pmatrix} 
               yh\left( x,y \right)\\ 
                -\frac{y h\left( x,y \right) }{x^{2} }  
         \end{pmatrix}  
         \]因此 \[
         \,\mathrm{d} _{\left( x,y \right) }F\left( Y \right)=  yh\left( x,y \right)\frac{\partial }{\partial u}-\frac{yh\left( x,y \right)  }{x^{2} }\frac{\partial }{\partial v}   
         \]
        \item 记 \(  \tilde{X}= d F\left( X \right)   \), \(  \tilde{Y}=  d F\left( Y \right)   \)  , 并记 \[
        \tilde{X}=  \tilde{X}^{u}\frac{\partial }{\partial u}+ \tilde{X}^{v}\frac{\partial }{\partial v},,\quad \tilde{Y}= \tilde{Y}^{u}\frac{\partial }{\partial u}+ \tilde{Y}^{v}\frac{\partial }{\partial v}
        \]则\[
        \left[ \tilde{X}, \tilde{Y} \right] ^{u}= \tilde{X}\left( \tilde{Y}^{u} \right)- \tilde{Y}\left( \tilde{X}^{u} \right)  
        \] \[
        \left[ \tilde{X},\tilde{Y} \right]^{v}= \tilde{X}\left( \tilde{Y} ^{v}\right)-\tilde{Y}\left( \tilde{X}^{v} \right)   
        \] \[
        \begin{aligned}
        \left( \tilde{X}\left( \tilde{Y}^{u} \right)  \right)\circ F&=  X\left( \tilde{Y}^{u}\circ F \right)= X\left( yh \right)  \\ 
         &= \left( f\left( x \right)\frac{\partial }{\partial x}+ g\left( y \right)\frac{\partial }{\partial y}   \right)   \left( yh \right)\\ 
          &= yf\frac{\partial h}{\partial x}+ gh+ yg\frac{\partial h}{\partial y} 
        \end{aligned}
        \] \[
        \begin{aligned}
        \left( \tilde{Y}\left( \tilde{X}^{u} \right)  \right)\circ F&= Y\left( \tilde{X}^{u}\circ F \right)   = Y\left( yf+ xg \right) \\ 
         &= \left( h\left( x,y \right)\frac{\partial }{\partial x}  \right)\left( yf\left( x \right)+ xg\left( y \right)   \right)\\ 
          &= yhf^{\prime} \left( x \right)+ hg\left( y \right)    
        \end{aligned}
        \]于是 \[
        \begin{aligned}
        \left( \left[ \tilde{X},\tilde{Y} \right]^{u}  \right)\circ F&=  \left( \tilde{X}\left( \tilde{Y}^{u} \right)  \right)\circ F-\left( \tilde{Y}\left( \tilde{X}^{u} \right)  \right)\circ F\\ 
         &= y\left( f\frac{\partial h}{\partial x}+ g\frac{\partial h}{\partial y}-hf^{\prime}  \right)     
        \end{aligned}
        \]接下来 \[
        \begin{aligned}
        \left( \tilde{X}\left( \tilde{Y} ^{v}\right)  \right)\circ F&= X\left( \tilde{Y}^{v}\circ F \right)= X\left( -\frac{y }{x^{2} }h  \right)\\ 
         &=  f\frac{\partial }{\partial x}\left( -\frac{y }{x^{2} }h  \right)+ g\frac{\partial }{\partial y}\left( -\frac{y }{x^{2} }h  \right)\\ 
          &= \frac{2yfh }{x^{3} }-\frac{yf }{x^{2} }\frac{\partial h}{\partial x}- \frac{gh }{x^{2} }- \frac{yg }{x^{2} }\frac{\partial h}{\partial y}          
        \end{aligned}
        \] \[
        \begin{aligned}
        \left( \tilde{Y}\left( \tilde{X}^{v} \right)  \right)\circ F&=  Y\left( \tilde{X}^{v}\circ F \right)   = Y\left( -\frac{y }{x^{2} }f+ \frac{g }{x }   \right)\\ 
         &= h\frac{\partial }{\partial x}\left( -\frac{y }{x^{2} }f+ \frac{g }{x }   \right)\\ 
        &= \frac{2yhf }{x^{3} }-\frac{yhf^{\prime}  }{x^{2} }-\frac{hg }{x^{2} }     
        \end{aligned}
        \]于是 \[
        \begin{aligned}
        \left( \left[ \tilde{X},\tilde{Y} \right]^{v}  \right)\circ F&= \left( \tilde{X}\left( \tilde{Y}^{v} \right)  \right)\circ F-\left( \tilde{Y}\left( \tilde{X}^{v} \right)  \right)\circ F\\ 
         &= -\frac{y }{x^{2} }\left( f\frac{\partial h}{\partial x}+ g\frac{\partial h}{\partial y}-hf^{\prime}  \right)      
        \end{aligned}
        \]令 \(  G\left( x,y \right)= f\frac{\partial h}{\partial x}+ g\frac{\partial h}{\partial y}-hf^{\prime}    \) , 则 \[
        \left[ \,\mathrm{d} F\left( X \right),\,\mathrm{d} F\left( Y \right)   \right]= y G\left( x,y \right)\frac{\partial }{\partial u}- \frac{y }{x^{2} }G\left( x,y \right)    \frac{\partial }{\partial v}
        \]

        另一方面,  \[
       \begin{aligned}
        \left[ X,Y \right]&= \left( X\left( h \right)-Y\left( f \right)   \right)\frac{\partial }{\partial x}+ \left( X\left( 0 \right)-Y\left( g \right)   \right)\frac{\partial }{\partial y}\\ 
         &= \left( \left( f\frac{\partial }{\partial x}+ g\frac{\partial }{\partial y} \right)  \left( h \right)- \left( h\frac{\partial }{\partial x} \right)f\left( x \right)   \right)\frac{\partial }{\partial x}- \left( \left( h\frac{\partial }{\partial x} \right)\left( g\left( y \right)  \right)   \right)\\ 
          &=  \left( f\frac{\partial h}{\partial x}+ g\frac{\partial h}{\partial y}-hf^{\prime}  \right)\frac{\partial }{\partial x}\\ 
           &= G\left( x,y \right)\frac{\partial }{\partial x}       
       \end{aligned} 
        \] \[
        \,\mathrm{d} F\left( \left[ X,Y \right]  \right)=  y G\left( x,y \right)\frac{\partial }{\partial u}- \frac{y }{x^{2} }G\left( x,y \right)\frac{\partial }{\partial v}    
        \]故 \[
        \left[ \,\mathrm{d} F\left( X \right),\,\mathrm{d} F\left( Y \right)   \right]= \,\mathrm{d} F\left( \left[ X,Y \right]  \right)  
        \]
        这是一个普遍规律. 为此, 记 \(  \tilde{X}= d F\left( X \right)   \), \(  \tilde{Y}=  d F\left( Y \right)   \). 则 对于任意的 \(  g \in C^{\infty}\left( M \right)   \), \[
        \begin{aligned}
        \left( \left[ \tilde{X}, \tilde{Y} \right]\left( g \right)   \right)\circ F&= \left( \tilde{X}\left( \tilde{Y}\left( g \right)  \right)  \right)\circ F-\left( \tilde{Y}\left( \tilde{X}\left( g\right)  \right)\right)    \circ  F\\ 
         &=  X\left( \left( \tilde{Y}\left( g \right)  \right)\circ F  \right) - Y\left( \left( \tilde{X}\left( g \right)  \right)\circ F  \right)\\ 
          &= X\left( Y\left( g\circ F \right)  \right)- Y\left( X\left( g\circ F \right)  \right)    \\ 
           &= \left[ X,Y \right]\left( g\circ F \right)  
        \end{aligned}
        \]   另一方面,  \[
        \left( \,\mathrm{d} F\left( \left[ X,Y \right]  \right)\left( g \right)   \right)\circ F= \left[ X,Y \right]\left( g\circ F \right)   
        \]二者相等, 故 \[
        \left[ \tilde{X}, \tilde{Y} \right]= \,\mathrm{d} F\left( \left[ X,Y \right]  \right)  
        \]当 \(  F  \)是同胚时成立. 
    \end{enumerate}
    
\end{proof}

\end{document}